\section{Exact Fourier Identity and Non-Approximation}
\label{sec:fourier-identity}

Orthogonality identifies the metric of $B_N$, but not its frequencies.  Equality
with the DFT follows from evaluation, independently of the round-trip theorem.

\begin{definition}[Packed real DFT]
For
\begin{equation}
 X_k=\sum_{n=0}^{N-1}x_n e^{-2\pi i kn/N},
 \label{eq:dft-definition}
\end{equation}
define $\DFT_N^{\R}:\R^N\to\R^N$ by
\begin{equation}
\begin{aligned}
 \DFT_N^{\R}x=(&X_0,\sqrt2\Re X_1,-\sqrt2\Im X_1,\ldots,\\
 &\sqrt2\Re X_{N/2-1},-\sqrt2\Im X_{N/2-1},X_{N/2})^{\trans}.
\end{aligned}
 \label{eq:packed-dft}
\end{equation}
\end{definition}

\begin{definition}[Normalized Bruun map]
Apply $\red_N$ in \eqref{eq:global-reduction}.  Express each quadratic residue
uniquely as $a_k+b_ke_{2\pi k/N}$, retain the scalar residues at $z=1,-1$,
multiply every quadratic pair $(a_k,b_k)$ by $\sqrt2$, and list these real
coordinates in native factor-tree order.  The resulting real linear map is
\begin{equation}
 B_N:\R^N\longrightarrow\R^N.
 \label{eq:BN-definition}
\end{equation}
The terminal pair weight is part of $B_N$.
\end{definition}

\begin{lemma}[Normalized reduction diagram]
Let $\mathcal A_t$ be the product of quotient algebras at a cut through the
Bruun tree, retaining completed scalar and quadratic leaves, and let
$\rho_t:\mathcal A_t\to\mathcal A_{t+1}$ be the simultaneous remainder map at
the active nodes.  Let $C_t$ list active polynomial coefficients, completed
scalars, and completed quadratic coefficients in the ordered basis
$(1,e_\theta)$.  Define
\begin{equation}
 S_t=C_{t+1}\rho_tC_t^{-1}.
 \label{eq:normalized-rank}
\end{equation}
Then the diagram $C_{t+1}\rho_t=S_tC_t$ commutes and
\begin{equation}
 B_N=W_NC_T\rho_{T-1}\cdots\rho_0C_0^{-1},
 \label{eq:BN-stage-product}
\end{equation}
where $W_N$ is $\sqrt2 I_2$ on each completed quadratic pair and identity on
the scalar leaves.  In a pair-oriented cut schedule, every generic block of
$S_t$ is the cell $T_\theta$, every ridge block is $H_2$, and completed leaves
are carry blocks.
\end{lemma}
\begin{proof}
The commuting identity is the definition of a change of coordinates on the
exact remainder map.  Multiplying the identities telescopes to the complete
CRT reduction, and $W_NC_T$ is exactly Definition~\ref{eq:BN-definition}.
For a binomial ridge, coefficient reduction is the add-subtract map $H_2$.
For a conjugate pair, evaluation of the two child representatives gives
\eqref{eq:complex-pair-merge}; that equation gives
the displayed real cell.  A factor already terminal at a cut is unchanged and
therefore carried.  The terminal $W_N$ supplies the declared pair packing.
\end{proof}

\begin{lemma}[Evaluation commutes with reduction]
If $q(\omega)=0$ and $f\equiv r\pmod q$, then $f(\omega)=r(\omega)$.
Consequently, the value at a terminal root is unchanged by every remainder map
on the path from the root polynomial to its leaf.
\end{lemma}
\begin{proof}
Write $f-r=qh$ and evaluate at $\omega$.  The path statement follows by
composition.
\end{proof}

Let $P_N$ be the signed permutation matrix that sends native factor-tree order to the
order in \eqref{eq:packed-dft}.  Its address law is made explicit in Section
\ref{sec:orders}; all scaling is already part of $B_N$, and
$P_N^{\trans}P_N=I$.

\begin{theorem}[Exact Fourier identity]
For every $N=2^L$ and every $x\in\R^N$,
\begin{equation}
 \boxed{\;P_NB_Nx=\DFT_N^{\R}x\;}.
 \label{eq:exact-fourier-identity}
\end{equation}
The equality holds as a linear identity over $\R$.  Its matrix coefficients lie
in the algebraic field $\K_N$ of \eqref{eq:coefficient-field}.
\end{theorem}
\begin{proof}
Associate $x$ with the polynomial in \eqref{eq:intro-polynomial}.  The forward
tree computes its exact residue at every terminal factor.  The scalar leaves
$z-1$ and $z+1$ evaluate to $X_0$ and $X_{N/2}$.  At $q_k$, write the normalized
residue as $a_k+b_ke_{2\pi k/N}$.  By the preceding lemma and the normalized
chart proposition,
\[
 X_k=x(e^{-2\pi ik/N})=a_k-ib_k.
\]
Conjugacy gives $X_{N-k}=\overline{X_k}$.  The terminal pair scaling and $P_N$
place $(\sqrt2a_k,\sqrt2b_k)$ in exactly the coordinates of
\eqref{eq:packed-dft}.  Every packed coordinate therefore agrees.
\end{proof}

\begin{corollary}[Global normalization, inverse, and Parseval identities]
\begin{equation}
\begin{aligned}
 B_N^{\trans}B_N&=NI,\\
 (\DFT_N^{\R})^{-1}&={1\over N}(\DFT_N^{\R})^{\trans},\\
 \|\DFT_N^{\R}x\|_2^2&=N\|x\|_2^2.
\end{aligned}
 \label{eq:inverse-parseval}
\end{equation}
Thus $B_N/\sqrt N$ is orthogonal, every singular value is $\sqrt N$, and the
condition number is one.  In a normalized cell schedule, the local inverse
halves and boundary weights distribute the inverse measure; no separate
terminal normalization pass is required.
\end{corollary}

\begin{corollary}[Fourier modes and convolution]
Real cosine and sine modes map to their corresponding single conjugate-pair
leaf, while the constant and alternating modes map to the scalar leaves.  For
circular convolution, multiplication of the complex numbers represented by
the terminal pairs gives the ordinary convolution theorem after removing the
boundary $\sqrt2$ weight from each input pair and restoring it on output.
\end{corollary}

Equation \eqref{eq:exact-fourier-identity} is the relevant non-approximation
statement.  $B_N$ is an exact factorization of the DFT linear map.  A
finite-precision program evaluates this factorization with rounded coefficients
and arithmetic; those residuals do not define the transform.
